% !PRINTED 115
\weberpage{115}
%<<CONT

\begin{equation}
f(x)=x^n+a_1x^{n-1}+a_2x^{n-2}+\cdots+a_n
\tag{1}
\end{equation}
the coefficient of $x^n$ equal to $1$, which in general is obtained by division by the coefficient of $x^n$. Then by the simple substitution
\begin{equation}
x=y-\frac{a_1}{n},
\tag{2}
\end{equation}
we can transform $f(x)$ into another function $\varphi(y)$ of the $n$th degree in which the second term $y^{n-1}$ does not occur; for, by the binomial theorem,
\[
 x^n=y^n-a_1y^{n-1}+\frac{n-1}{2n}a_1^2y^{n-2}+\cdots,
\]
\[
 x^{n-1}=y^{n-1}-\frac{n-1}{n}a_1y^{n-2}+\cdots,
 \qquad x^{n-2}=y^{n-2}-\cdots,
\]
and consequently
\[
 f(x)=y^n+\left(a_2-\frac{n-1}{2n}a_1^2\right)y^{n-2}+\cdots .
\]

If we actually carry out the transformation for the cases $n=2,3,4$, we obtain
\begin{equation}
\begin{array}{rcl}
n=2, & \varphi(y)&=y^2+a,\\[0.25em]
n=3, & \varphi(y)&=y^3+ay+b,\\[0.25em]
n=4, & \varphi(y)&=y^4+ay^2+by+c,
\end{array}
\tag{3}
\end{equation}
provided that we put
\begin{equation}
\begin{array}{rcll}
\text{for }n=2:&& a&=-\dfrac{a_1^2}{4}+a_2,\\[0.75em]
\text{for }n=3:&& a&=-\dfrac{a_1^2}{3}+a_2,\\[0.75em]
&& b&=\dfrac{2a_1^3}{27}-\dfrac{a_1a_2}{3}+a_3,\\[0.75em]
\text{for }n=4:&& a&=-\dfrac{3a_1^2}{8}+a_2,\\[0.75em]
&& b&=\dfrac{a_1^3}{8}-\dfrac{a_1a_2}{2}+a_3,\\[0.75em]
&& c&=-\dfrac{3a_1^4}{256}+\dfrac{a_1^2a_2}{16}-\dfrac{a_1a_3}{4}+a_4.
\end{array}
\tag{4}
\end{equation}
Thus the solution of the equation $f(x)=0$ is reduced by (2) to the solution of $\varphi(y)=0$. This first gives, for $n=2$,
%CONT>>
